% from ega1-6-fr.tex
\oldpage[I]{147}
facteurs $X$, $Y$ est \emph{de type fini} sur $S$ et
\emph{l'autre noethérien} (resp. \emph{localement noethérien}), alors
$X\times_S Y$ est \emph{noethérien} (resp. \emph{localement noethérien}).

\begin{corollary}[6.3.9]
\label{I.6.3.9-fr}
Soit $X$ un préschéma de type fini sur un préschéma localement noethérien $S$.
Alors tout $S$-morphisme $f:X\to Y$ est de type fini.
\end{corollary}

En effet, on peut supposer $S$ noethérien~; si $\varphi:X\to S$,
$\psi:Y\to S$ sont les morphismes structuraux, on a $\varphi=\psi\circ f$, et
$X$ est noethérien en vertu de (\hyperref[I.6.3.7-fr]{6.3.7})~; $f$ est donc
de type fini en vertu de (\hyperref[I.6.3.6-fr]{6.3.6}).

\begin{proposition}[6.3.10]
\label{I.6.3.10-fr}
Soit $f:X\to Y$ un morphisme de type fini. Pour que $f$ soit surjectif, il faut
et il suffit que, pour tout corps algébriquement clos $\Omega$, l'application
$X(\Omega)\to Y(\Omega)$ correspondant à $f$
(\hyperref[I.3.4.1-fr]{3.4.1}) soit surjective.
\end{proposition}

La condition est suffisante, comme on le voit en considérant pour tout
$y\in Y$, une extension algébriquement close $\Omega$ de $k(y)$, et le
diagramme commutatif
\[
  \xymatrix{
    X\ar[dd]_f\\
    & \operatorname{Spec}(\Omega)\ar[ul]\ar[dl]\\
    Y
  }
\]
(\emph{cf.} (\hyperref[I.3.5.3-fr]{3.5.3})). Inversement, supposons $f$
surjective, et soit $g:\{\xi\}=\operatorname{Spec}(\Omega)\to Y$ un morphisme,
$\Omega$ étant un corps algébriquement clos. Si on considère le diagramme
\[
  \xymatrix{
    X\ar[d]_f &
    X_{(\Omega)}\ar[l]\ar[d]^{f_{(\Omega)}}\\
    Y &
    \operatorname{Spec}(\Omega)\ar[l]
  }
\]
il suffit donc de montrer qu'il existe dans $X_{(\Omega)}$ un
\emph{point rationnel sur $\Omega$}
(\hyperref[I.3.3.14-fr]{3.3.14}, \hyperref[I.3.4.3-fr]{3.4.3} et
\hyperref[I.3.4.4-fr]{3.4.4}). Comme $f$ est surjectif, $X_{(\Omega)}$ n'est
pas vide (\hyperref[I.3.5.10-fr]{3.5.10}) et comme $f$ est de type fini, il en
est de même de $f_{(\Omega)}$ (\hyperref[I.6.3.4-fr]{6.3.4, (iv)})~; donc
$X_{(\Omega)}$ contient un ouvert affine non vide $Z$ tel que $A(Z)$ soit une
algèbre de type fini non nulle sur $\Omega$. En vertu du th. des zéros de
Hilbert \cite{I-21}, il existe un $\Omega$-homomorphisme $A(Z)\to\Omega$, donc
une section de $X_{(\Omega)}$ au-dessus de $\operatorname{Spec}(\Omega)$, ce
qui démontre la proposition.

\subsection{Préschémas algébriques.}
\label{subsection:I.6.4-fr}

\begin{definition}[6.4.1]
\label{I.6.4.1-fr}
Étant donné un corps $K$, on appelle \emph{$K$-préschéma algébrique} un
préschéma $X$ de type fini sur $K$~; $K$ est appelé le \emph{corps de base}
de $X$. Si en outre $X$ est un schéma (ou, ce qui revient au même
(\hyperref[I.5.5.8-fr]{5.5.8}), si $X$ est un $K$-schéma), on dit aussi que
$X$ est un \emph{$K$-schéma algébrique}.
\end{definition}

Tout $K$-préschéma algébrique est \emph{noethérien}
(\hyperref[I.6.3.7-fr]{6.3.7}).

\begin{proposition}[6.4.2]
\label{I.6.4.2-fr}
Soit $X$ un $K$-préschéma algébrique. Pour qu'un point $x\in X$ soit fermé,
il faut et il suffit que $k(x)$ soit une extension algébrique de $K$, de degré
fini.
\end{proposition}

On peut supposer $X$ affine, l'anneau $A$ de $X$ étant une $K$-algèbre de type
fini. En effet, les ouverts affines $U$ de $X$ tels que $A(U)$ soit une
$K$-algèbre de type fini forment un recouvrement de $X$
(\hyperref[I.6.3.1-fr]{6.3.1}). Les points fermés de $X$ sont alors les points
tels que $\mathfrak{j}_x$ soit
